\documentclass{article}
\usepackage[french]{babel}

\usepackage{mathrsfs}
\usepackage{mathtools}
\usepackage{amsmath}
\usepackage{amsthm} 
\usepackage{amsfonts}
\usepackage{aliascnt}
\usepackage{dsfont}
\usepackage{enumitem}
\usepackage{chngcntr}
\usepackage[dvipsnames]{xcolor}
\usepackage{xparse}
\usepackage{bm}
\usepackage{bbm}
\usepackage{listings}
\usepackage[most]{tcolorbox}

% newcomands ==============================
\newcommand{\F}{\mathscr{F}}
\newcommand{\N}{\mathscr{N}}
\newcommand{\carE}{\mathscr{E}}
\renewcommand{\P}{\mathds{P}}
\newcommand{\R}{\mathbb{R}}
\newcommand{\Q}{\mathbb{Q}}
\newcommand{\E}{\mathbb{E}}
\newcommand{\1}{\mathbbm{1}}
\newcommand{\sign}{\text{sign}}
\newcommand{\supp}{\text{supp}}
\newcommand{\refocc}{\hyperref[eq:occupation]{$(P_{t,\varphi})$}}
\newcommand{\refcontocc}{\hyperref[proof:continuite_occ]{\textbf{Continuité de \refocc\; en $t$}}}
\newcommand{\reffin}{\hyperref[proof:finfinfin]{$(\star)$}}

\renewcommand{\d}{\mathrm{d}}

\RenewDocumentCommand{\footnote}{m m}{%
    \textit{#1}\footnotemark\footnotetext{#2}%
  }

\newcounter{questionCounter}
\setcounter{questionCounter}{0}
\newcommand{\Question}[1]{%
    \stepcounter{questionCounter} % Incrémente le compteur
    \bigskip\noindent {\color{BlueViolet}\textbf{ Question \thequestionCounter} : #1}
    \par % Saut de ligne
  }

% Style code   ============================
\lstdefinestyle{customc}{
  language=Python,  
  breaklines=false,
  backgroundcolor=\color{gray!10},
  basicstyle=\footnotesize\ttfamily,
  keywordstyle=\bfseries\color{green!40!black},
  commentstyle=\itshape\color{purple!40!black},
  identifierstyle=\color{blue},
  stringstyle=\color{orange},
  numberstyle=\tiny\color{mygray},
}

\lstset{escapechar=@,style=customc}
% Mise en page ============================

\title{Options sur Best-of}
\author{Hugo De Brito, Quentin Beaurez et Héléna Benet Burgaud}
% =========================================
\begin{document}
\maketitle

\section{Introdution}
Les options \textbf{Best-of} sont des produits financiers dont la valeur à l'échéance dépend de la performance du \textbf{meilleur actif} parmis un panier de sous-jacents.


On considerera d'abord un \footnote{contrat forward}{Le {contrat forward} est un accord d'achat ou de vente d'un actif à un prix et une date future précisée dans un contrat.} sur une option Best-of puis on considerera le cas d'un \footnote{put}{Le \textit{put} désigne une option de vente}. Cette option est dépendante de la loi jointe des actifs, on étudiera en particulier cet aspect. Pour cela on utilisera la méthode de valorisation par téchnique de \textbf{Monte Carlo} dans le cadre du modèle de Black \& Scholes.

\subsection{Modèle}

On considère le modèle de Black \& Scholes de dimension 3. Soient $\bm S_t = (S^1_t, S^2_t, S^3_t)$ le vecteur des actifs risqués au temps $t$
et $\bm S_0 = (s^1_0, s^2_0, s^3_0)$ le vecteur de valeurs initiales des actifs supposé déterministe. Soit $\bm W_t = (W^1_t, W^2_t, W^3_t)$ un moment brownien sous probabilité \footnote{risque neutre}{...} de matrice de covariation $\bm \Gamma$. Soit $r$ le taux d'interêts supposé constant et soit $\bm \sigma = (\sigma_1, \sigma_2, \sigma_3)$ le vecteur de volatilités.

$S^i_t$ est solution de l'EDS :

\begin{equation}
  \begin{cases}\label{EDS:BlackScholes}
    \d S^i_t &= S^i_t (r \d t + \sigma_i \d W^i_t)\\
    S^i_0 &= s^i_0 > 0
  \end{cases}
\end{equation}

avec
\begin{equation*}
\bm \Gamma = \begin{pmatrix}
          1 & \rho & \rho \\
          \rho & 1 & \rho \\
          \rho & \rho & 1 \\
        \end{pmatrix}
      \end{equation*}

et $\rho \in (-1/2, 1]$.

%%% QUESTION 1 %%%%%%%%%%%%%%%%%%%%%%%%%%%%%%%%%%%%%%%%%%%%%%%%%%%%%%

\Question{Calculer $\bm S_t$}
La solution de \eqref{EDS:BlackScholes} est facilement calculable. On peut supposer sans perte de généralité que $S^i_t > 0$ ps; ceci nous permet d'écrire
$$\frac{\d S^i_t}{S^i_t} = r \d t + \sigma_i \d W^i_t $$

On reconnait la dérivée de $\log(S^i_t)$. En notant que $\log(\cdot)$ est $\mathcal C^2$ sur $\R^*_+$

\begin{align*}
  \log(S^i_t) &\overset{\text{Itō}}{=} \log(s^i_t) + \int_0^t \dfrac{1}{S^i_s} \d S^i_s + \dfrac{1}{2}\int_0^t -\dfrac{1}{(S^i_s)^2} \d [S^i]_s \quad p.s.
  \\
              &= \log(s^i_t) + \int_0^t r - \dfrac{\sigma_i^2}{2} \d s + \int_0^t\sigma_i \d W^i_s \quad p.s.\\
              &= \log(s^i_t) + \left(r - \dfrac{\sigma_i^2}{2}\right)t + \sigma_i W^i_t \quad p.s.
\end{align*}

Donc 
\begin{equation}\label{Sol:BlackScholes}
  S^i_t = \exp\left(\log( S^i_t)\right) = s^i_0\exp\left(\left(r - \dfrac{\sigma_i^2}{2}\right)t + \sigma_i W^i_t\right) \quad p.s.\\ 
\end{equation}

Ainsi $\bm S_t = (S^1_t, S^2_t, S^3_t)$ où les $S^i_t$ sont de la forme \eqref{Sol:BlackScholes}.

\section{Forward sur Best-of}

Un contrat forward sur Best Of confère à son possesseur l'obligation d'acheter l'actif le plus cher parmi plusieurs actifs prédéfinis à un prix fixé à l'avance. Le prix de ce forward sur Best Of de 3 actifs est donc dans notre modèle

$$P^1 \ = \ e^{-rT} \E \left( \max_i(S^i_T) - K \right)$$

sous probabilité Risque Neutre.

%%% QUESTION 2 %%%%%%%%%%%%%%%%%%%%%%%%%%%%%%%%%%%%%%%%%%%%%%%%%%%%%%

\Question{Montrer que le prix vérifie $$ P^1 \geq \max_{i} (s^i_0) - e^{-rT}K$$}

Dans quels cas simples en fonction de $\sigma$ et $\rho$ a-t-on égalité (expliquer) ?
Soient $T>0$ et $K\in \R$ une constante déterministe. On prend  $t \in [0, T]$. 
On se donne $\tilde S^i_t = e^{-rt} S^i_t$ le prix actualisé.
\begin{align*}
  \tilde S^i_t &\overset{\text{Itō}}{=} s^i_0 + \int_0^t -r e^{-rs} S^i_s \d s + \int_0^t e^{-rs} \d S^i_s,\\
               &= s^i_0 +  \int_0^t \sigma_i e^{-rs} S^i_s \d W^i_s\\
  &= s^i_0 + \int_0^t \sigma_i \exp\left(\sigma_i W^i_s - \frac{1}{2}\sigma_i^2 s\right) \d W^i_s
\end{align*}

\begin{equation}
  P^1 = \E\left(\max_i(\tilde S^i_T) - K\right) =  \E\left(\max_i(\tilde S^i_T)\right) - K \,e^{-rT}
\end{equation}

Notons que la fonction $\max(\cdot)$ est une fonction convexe. Par Jensen, on a :

\begin{align*}
  P^1 &\geq \max_i \E\left(\tilde S^i_T\right) - K\,e^{-rT}\\
      &=  \max_i\left( \E\left(s^i_0 + \int_0^T \sigma_i \exp(\sigma_i W^i_s - \frac{1}{2} \sigma_i^2 s) \d W^i_s\right)\right) - K\,e^{-rT}\\
      &= \max_i\Bigg(s^i_0 +  \E\Big(\int_0^T \sigma_i \exp(\sigma_i W^i_s - \frac{1}{2} \sigma_i^2 s) \d W^i_s\Big)\Bigg) - K\,e^{-rT}\\
  \intertext{Notons que l'intégrale stochastique $\int_0^T \sigma_i \exp(\sigma_i W^i_s - \frac{1}{2} \sigma_i^2 s) \d W^i_s$ suit une loi normale centrée.}
  &\geq \max_i(s^i_0) - K\,e^{-rT}
\end{align*}

%%%%% NOTE 
Notons de plus que $(\tilde S^i_t)_{t\in[0,T]}$ est une martingale. En effet, comme on peut exprimer  $\tilde S^i_t$ comme une intégrale stochastique : il suffit de vérifier que :
\begin{equation}\label{Q2:condMtmart}
  \E\left(\int_0^t \sigma_i^2 \exp\left(2 \sigma_i W^i_s - \sigma_i^2 s\right) \d s\right) < \infty \;p.s.
  \end{equation}

  Soit $M^i_t = \exp\left(2 \sigma_i W^i_s - \sigma_i^2 s\right)$. La condition \eqref{Q2:condMtmart} est vérifiée puisque $(M_t)_{t \in [0,T]}$ est une martingale, Ainsi, 

  \begin{alignat}{2}
    \text{\eqref{Q2:condMtmart}} &&\quad\Leftrightarrow\quad \E\left(\int_0^T \sigma_i^2 M_s \d s\right) &\quad=\quad \int_0^T \sigma_i^2\E(M^i_s) \d s \\
    &&\quad=\quad \int_0^T \sigma_i^2\E(M^i_0) \d s &\quad=\quad \sigma_i^2\E(M^i_0) T\quad<\quad \infty \;p.s.
    \end{alignat}

  Montrons que $(M^i_t)_{t \in[0,T]}$ est une martingale. Tout d'abord, notons que on peut réecrire $M_t$ comme suit :

$$M^i_t = \exp\left( - \int_0^t (-\sqrt{2}\sigma_i) \d W^i_s - \frac{1}{2}\int_0^t (-\sqrt{2}\sigma_i)^2 \d s\right).$$
Par le critère de Novikov, $M^i_t$ est une martingale si $\E\left(\exp\left(\int_0^T \frac{1}{2} (-\sqrt{2}\sigma_i)^2 \d s\right)\right) < \infty \; p.s.$ Or, on a :
\begin{equation*}
  \E\left(\exp\left(\int_0^T \frac{1}{2} (-\sqrt{2}\sigma_i)^2 \d s\right)\right) = \E \left(\exp\left(\sigma_i^2 T\right)\right) = \exp\left(\sigma_i^2 T\right)< \infty \; p.s.
\end{equation*}


Ainsi, par la propriété de martingale $\E(M^i_t) = \E(M^i_0) = 1$. Ceci conclut que le prix actualisé $\tilde S^i_t$ est bien une martingale.

%%%

%%% QUESTION 3 %%%%%%%%%%%%%%%%%%%%%%%%%%%%%%%%%%%%%%%%%%%%%%%%%%%%%%

\Question{Proposer une méthode pour simuler $W(T)$. Implémenter le calcul d'un estimateur de $P^1$ par une méthode de Monte Carlo classique. On choisira pour les applications $\sigma_i = 0.3$, $S_{i,0} = 1$ $\forall i$ (sauf pour Q11), $\rho = 0.3$ (sauf pour Q6-Q10), $K = 1$ (sauf pour Q10), $r = 0.02$ et $T = 1.5$ ans. Les volatilités et le taux sont annuels.\label{Q4}}


Pour simpuler $\bm W(T)$ il suffit de simuler en premier les trajectoires des mouvements browniens $W^1$, $W^2$ et $W^3$ puis de les corréler à l'aide de la décomposition de Chloesky de la matrice de covariation $\bm \Gamma$.

%Pour simuler les trajectoires du mouvement brownien on note que : \\$ W_t \overset{\text{Loi}}{=} \sqrt{t}\;N$ où $N \sim \mathcal N(0,1)$.
%
%Ainsi il faut en premier simuler la loi normale $\mathcal N(0,1)$ grâce à l'algorithme de \textbf{Box Muller}.
%
%%\lstinputlisting[language=Python, firstline=12, lastline=25]{BestOf.py}
%{\color{violet} \textbf{code}}
%
%
%Dès qu'on a une simulation de la loi normale on peut modéliser des trajectoires de mouvement browniens.
%
%%\lstinputlisting[language=Python, firstline=34, lastline=55]{BestOf.py}
%{\color{violet} \textbf{code}}
%
%
%Finalement, pour corréler les mouvements browniens, on utilise la fonction \texttt{np.linalg.cholesky} de python. 
%
%%\lstinputlisting[language=Python, firstline=78, lastline=91]{BestOf.py}
%{\color{violet} \textbf{code}}

%%% QUESTION 4 %%%%%%%%%%%%%%%%%%%%%%%%%%%%%%%%%%%%%%%%%%%%%%%%%%%%%%

\Question{Implémenter le calcul d'un estimateur de $P^1$ par une méthode de Monte Carlo utilisant une technique de réduction de variance par variables antithétiques.}

%%% QUESTION 5 %%%%%%%%%%%%%%%%%%%%%%%%%%%%%%%%%%%%%%%%%%%%%%%%%%%%%%

\Question{Calculer les variances empiriques (des estimateurs) obtenues en fonction du nombre de trajectoires pour les estimations (avec et sans réduction de variance) Monte Carlo de $P^1$. Tracer sur le même graphique les estimateurs de $P^1$ en fonction du nombre de trajectoires ainsi que leurs intervalles de confiance asymptotiques à 90\%.}

%%% QUESTION 6 %%%%%%%%%%%%%%%%%%%%%%%%%%%%%%%%%%%%%%%%%%%%%%%%%%%%%%
\Question{Choisir un nombre de trajectoires pour lequel l'estimation est précise
(avec réduction de variance seulement) (préciser votre critère) et tracer sur le
même graphique l'estimateur de prix $P^1$ (avec réduction de variance) du forward
en fonction de $\rho$ pour $\rho in (0.5,1)$ (garder les autres paramètres). Commenter.}

\section{Put sur Best-of}
%%% QUESTION 7 %%%%%%%%%%%%%%%%%%%%%%%%%%%%%%%%%%%%%%%%%%%%%%%%%%%%%%

Un \textbf{put} est une option qui confère à son possesseur le droit, mais pas l'obligation, de vendre un actif à un prix fixé à l'avance. On parle d'une option \textbf{Européene} si la date à laquelle le droit sur l'option est exercée (i.e. que l'actif est soit acheté soit vendu) est la même que la date d'échéance du contrat.


Le prix d'un put sur l'actif $S^i$ à la date d'échéance $T$ est :

$$P^{E,i} = e^{-rT} \E\left( \left( K - S^i_T \right)_+\right)$$
sous probabilité risque neutre.

\Question{Calculer analytiquement et implémenter le prix de l’option $P^{E,i}$.}
{\color{violet}\textbf{Il manque }: Montrer que la $S^i_T$ suit une loi log-normale!!!}

Rappellons que la solution $S_T^i$ est de la forme $s^i_0 \, e^{(r-\frac{\sigma_i^2}{2})T + \sigma_i W_T^i}$

Ainsi on a :

\begin{align*}
  P^{E,i} &= e^{-rT} \E\left( \left( K - S^i_T \right)_+\right)\\
%          &= \E\left( \max(0, K\,e^{-rT} - s^i_0\, e^{-\frac{\sigma_i^2}{2}T + \sigma_i W_T^i} )\right)
            \intertext{Il est clair que dans le cas où $S_T^i \geq K$ alors le prix de l'option est égal à zéro. Ceci dit, on ne connait pas la valeur exacte de $S_T^i$, par contre on connait sa loi! En effet, on sait que $\log\left(\frac{S^i_T}{s^i_0}\right) \sim \N\left((r- \frac{\sigma_i^2}{2})\,T , \sigma_i^2 \,T\right)$. On peut dont exprimer l'esperance en fonction de la densité de $S_t^i$.}
          &= e^{-rT} \E\left( \max\Big(0 ,  K - s^i_0\,e^{\log\left(\frac{S_T^i}{s^i_0}\right)} \Big)\right)\\
          &= \int_{-\infty}^{+\infty} e^{-rT}\, \max\Big(0, K - s^i_0e^{\left(r - \frac{\sigma_i^2}{2}\right)\,T + \sigma_i \sqrt{T} x}\Big)p(x) \d x
            \intertext{où $p(x)$ est la densité de la loi $\N(0,1)$.}
\end{align*}

Remarquons que le $\max$ est supérieur à zéro quand $S_T^i\leq K $. c'est à dire :
\begin{align*}
        s_0^i\, e^{\left(r - \frac{\sigma_i^2}{2}\right)T + \sigma_i\sqrt{T}\,x}& \leq K \\
  \text{i.e.}\quad \sigma_i\sqrt{T}\,x &\leq \log\left(\frac{K}{s_0^i}\right) - \left(r - \frac{\sigma_i^2}{2}\right)T\\
  \text{i.e.}\quad x &\leq \frac{\log\left(\frac{K}{s_0^i}\right) - \left(r - \frac{\sigma_i^2}{2}\right)T}{\sigma_i\sqrt{T}} = - d_1\\
\end{align*}

Donc :

\begin{align*}
  P^{E,i} &= \int_{-\infty}^{- d_1} e^{-rT}\Big(K - s^i_0e^{\left(r - \frac{\sigma_i^2}{2}\right)\,T + \sigma_i \sqrt{T} x}\Big)p(x) \d x,\\
          &= \int_{-\infty}^{- d_1} K e^{-rT} \, p(x) \d x -  \int_{-\infty}^{- d_1}s^i_0e^{- \frac{\sigma_i^2}{2}\,T + \sigma_i \sqrt{T} x}p(x) \d x,
  \shortintertext{En notant que $K e^{-rT}$ est une constante par rapport à $x$, on a :}
  P^{E,i} &=  K e^{-rT}\, F(- d_1) - s^i_0\int_{-\infty}^{- d_1}\frac{1}{\sqrt{2 \pi}} e^{- \frac{\sigma_i^2}{2}\,T + \sigma_i \sqrt{T} x - \frac{x^2}{2}} \d x.
\shortintertext{où $F(\xi)$ est la fonction de répartition de la loi $\N(0,1)$.}
\intertext{Notons également que $- \frac{\sigma_i^2}{2}\,T + \sigma_i \sqrt{T} x - \frac{x^2}{2} = -\frac{(x - \sigma_i\sqrt{T})^2}{2}$. Ainsi on a :}
   P^{E,i} &=   K e^{-rT}\, F(- d_1) - s^i_0\int_{-\infty}^{- d_1}\frac{1}{\sqrt{2 \pi}}e^{-\frac{(x - \sigma_i\sqrt{T})^2}{2}} \d x,\\
\shortintertext{puis, en faisant le changement de variable $y = x - \sigma_i\sqrt{T}$ on voit que le deuxième terme est aussi une fonction de répartition d'une normale $\N(0,1)$. }
  P^{E,i}  &=   K e^{-rT} \, F(- d_1) - s^i_0\int_{-\infty}^{- d_1 - \sigma_i\sqrt{T}}\frac{1}{\sqrt{2 \pi}}e^{-\frac{y^2}{2}} \d y,\\
          &=  K e^{-rT}\, F(- d_1) - s_0^i\,F(- d_1 - \sigma_i\sqrt{T} ).\\
          &=  K e^{-rT}\, F(- d_1) - s_0^i\,F(-d_2 ).
\end{align*}


%%% QUESTION 8 %%%%%%%%%%%%%%%%%%%%%%%%%%%%%%%%%%%%%%%%%%%%%%%%%%%%%%

%Une option \textbf{put sur un Best-of} confère à son possesseur le droit mais pas l’obligation de \textbf{vendre} l’actif le \textbf{plus cher} parmi plusieurs actifs prédéfinis à un prix fixé à l’avance. Le prix de ce put sur Best Of de 3 actifs est donc dans notre modèle :
%
%
%$$P^2 = e^{-rT}\E\left(\left((K - \max_i(S^i_T)\right)_+\right)$$
%
%\Question{Déduire de la question precédente un majorant de $P^2$.}
%
%En notant que $(K - \max_i S_T^i)_+ >0$ que si pour tous les $i$, $S_T^i < K$. Ainsi on peut dire $(K - \max_i S_T^i)_+ \leq \max_i (K - S_T^i)_+$.
%
%Comme la constante $e^{-rT}$ est positive et par linéarité de l'esperance on a :
%$$P^2 = e^{-rT}\,\E\left(\left((K - \max_i(S^i_T)\right)_+\right) \leq e^{-rT}\,\E\left(\max_i  \left((K - S_T^i)_+\right) \right)$$
%
%Notons maintenant qu'on prend le maximum de v.a. positives. Ainsi, on a :
%$$\max_i(K = S_T^i)_+ \leq \sum_i (K-S_T^i)_+$$
%Donc :
%$$  P^2 \leq  e^{-rT}\,\E\left(\sum_i (K-S_T^i)_+\right) \leq \sum_i P^{E,i}.$$

Partons du fait que $$ \underset{i}{\text{moyenne}} S^i_T \leq \max_i S^i_T $$

Alors,

\begin{align*}
          & K - \max_i S^i_T  \leq K - \underset{i}{\text{moyenne}} S^i_T \\
\implies  & \left( K - \max_i S^i_T \right)+ \leq \left( K - \underset{i}{\text{moyenne}} S^i_T \right)+ \\
\implies  & \left( K - \max_i S^i_T \right)+ \leq \left( \underset{i}{\text{moyenne}} (K - S^i_T) \right)+ \leq \underset{i}{\text{moyenne}} \left\lbrace \left( K - S^i_T \right)_+ \right\rbrace
\end{align*}

D'où, 

$$ P^2 = e^{-rT} \E\left(\left((K - \max_i(S^i_T)\right)+\right) \leq   e^{-rT} \E \left( \dfrac{1}{n} \sum{i=1}^n  \left\lbrace \left( K - S^i_T \right)_+ \right\rbrace \right) \leq  e^{-rT} \, \underset{i}{\text{moyenne}} P^{E, i} $$
%%% QUESTION 9 %%%%%%%%%%%%%%%%%%%%%%%%%%%%%%%%%%%%%%%%%%%%%%%%%%%%%%
\Question{Implémenter le calcul d'un estimateur de $ P^2 $ par une méthode de Monte Carlo utilisant une technique de réduction de variance par variables antithétiques. Tracer sur le même graphique cet estimateur en fonction du nombre de trajectoires ainsi que ses intervalles de confiance asymptotiques à 90\%. Tracer également le même graphique dans le cas $ \rho = 1 $ et ajouter à celui-ci le précédent majorant. Commenter.}

%%% QUESTION 10 %%%%%%%%%%%%%%%%%%%%%%%%%%%%%%%%%%%%%%%%%%%%%%%%%%%%%%
\Question{Choisir un nombre de trajectoires pour lequel l'estimation est précise (préciser votre critère) et tracer sur le même graphique l'estimateur de prix $ P^2 $ en fonction de $ \rho $ pour $ \rho \in [-0.5, 1] $ (garder les autres paramètres) pour $ K = 1 $, $ K = 1.2 $ et $ K = 1.5 $. Commenter.}

%%% QUESTION 11 %%%%%%%%%%%%%%%%%%%%%%%%%%%%%%%%%%%%%%%%%%%%%%%%%%%%%% 
\Question{Tracer sur un même graphique le prix $ P^{E,1} $ et l'estimateur de prix $ P^2 $ en fonction de $ S_{1,0} $ pour $ S_{1,0} \in [0, 2] $ (garder les autres paramètres). Commenter.}

%%% QUESTION 12 %%%%%%%%%%%%%%%%%%%%%%%%%%%%%%%%%%%%%%%%%%%%%%%%%%%%%% 
\Question{Utiliser un put Européen sur $S_1$ comme variable de contrôle pour en déduire un nouvel estimateur du prix $P^2$. Illustrer graphiquement cette méthode.}

%%% QUESTION 13 %%%%%%%%%%%%%%%%%%%%%%%%%%%%%%%%%%%%%%%%%%%%%%%%%%%%%% 
Nous considérons cette fois un put sur Best Of sur $N$ sous-jacents. 
Les sous-jacents sont modélisés avec la même dynamique, avec
$$\Gamma = \left( \Gamma_{i,j} \right)_{i,j=1..N}, \quad \Gamma_{i,j} = \mathbf{1}_{i=j} + \rho \mathbf{1}_{i \neq j}$$

Le prix du Best Of est donc dans notre modèle
$$P^{2,N} = e^{-rT} \mathbb{E} \left[ \left( K - \max_{i=1..N}(S_T^i) \right)_+ \right]$$
sous probabilité Risque Neutre.

\Question{Choisir un nombre de trajectoires pour lequel l'estimation est précise et tracer sur un même graphique l'estimateur de prix $P^{2,N}$ en fonction de $N \in \{1, 3, 10, 25\}$ pour $\rho \in \{0, 0.5, 1\}$ (garder les autres paramètres) (on pourra utiliser une librairie externe pour correler les mouvements Browniens). Commenter.
}

\section{Call sur Best-Of}
Un call sur un Best-of est une option qui confère à son détenteur le droit (mais pas l'obligation) d'\textbf{acheter} un actif au choix parmi plusieurs actifs prédefinis à un prix fixé à l'avance.

Le prix de ce call sur Best Of de 3 actifs est donc dans notre modèle est :

$$P^3 = e^{-rT} \E \left( \left( \max_iS_T^i - K \right)_+ \right)$$

sous probabilité risque neutre.

%%% QUESTION 14 %%%%%%%%%%%%%%%%%%%%%%%%%%%%%%%%%%%%%%%%%%%%%%%%%%%%%% 

\Question{Quel lien existe-t-il entre le forward, le call, et le put sur Best of?}
Rappelons d'abord les prix d'un Best-of sur un contrat forward $P^1$, sur un put $P^2$ et sur un call $P^3$ :

\begin{align*}
  P^1 &= e^{-rT} \E \left( \max_i(S_T^i) - K \right)\\
  P^2 &=  e^{-rT} \E \left( \left( K - \max_i (S_T^i) \right)_+ \right)\\
  P^3 &= e^{-rT} \E \left( \left( \max_i(S_T^i)) - K \right)_+ \right)
\end{align*}

Notons de plus que pour deux réels $a, b \in \mathbb R$, $(a-b)_+ - (b-a)_+ = a-b \quad(\star)$. Ainsi, par linéaritéd de l'esperance, on a :

\begin{align}
  P^3 - P^2 &= e^{-rT} \E \left (\left( \max_i(S_T^i) - K \right)_+ - \left( K - \max_i(S_T^i)\right)_+  \right) \notag\\
            &= e^{-rT} \E\left( \max_i S_T^i - K \right)\notag\\
            &= P^1
\end{align}

Cette relation est connue en fincance sous le nom de \textbf{parité call-put}.


%%% QUESTION 15 %%%%%%%%%%%%%%%%%%%%%%%%%%%%%%%%%%%%%%%%%%%%%%%%%%%%%% 

\Question{ Calculer analytiquement le prix du forward sur Best-of sur 2 actifs
dans ce modèle}

Par la question précédente, on a que le prix d'un forward $P^1$ sur un Best-of satisfait la relation $P^1 = P^3 - P^2$ où $P^2$ est le prix d'un put sur un Best-of et $P^3$ est le prix d'un call sur un Best-of.

Notons que on cherche à calculer le prix du forwar sur un Best-of sur \textbf{deux} actifs. Donc le $\max_i S_T^i = \max(S_T^1, S_T^2) = S_T^2\max\left(\frac{S_T^1}{S_T^2},1\right)$.

Ainsi on a :

\begin{align*}
  P^1 &= P^3 - P^2\\
      &\overset{\text{Lin.}}{=} e^{-rT}\E\left( (\max(S_T^1, S_T^2)- K)_+ - (K -\max(S_T^1, S_T^2))_+\right)\\
      &\overset{(\star)}{=} e^{-rT}\E\left( \max(S_T^1, S_T^2)- K \right)
%  \shortintertext{Comme $S_T^2$ est non nul p.s. on peut écrire :}
%      &= e^{-rT}\E\left(S_T^2 \max\left(\frac{S_1^1}{S_T^2},1\right)\right) - e^{-rT}K 
  \shortintertext{Puis, en utilisant la défintion trajectorelle des $S_T^i$ :}
 P^1&=  e^{-rT} \E \left(\max\left(s_0^1 \exp\left((r- \frac{\sigma_1^2}{2})\,T + \sigma_1W_t^1\right), s_0^2 \exp\left((r- \frac{\sigma_2^2}{2})\,T + \sigma_2W_t^2\right) \right)\right)\\
      &\qquad- e^{-rT}K,\\
    &= \E\left(\max\left(s_0^1 \exp\left(- \frac{\sigma_1^2}{2}\,T + \sigma_1W_t^1\right), s_0^2 \exp\left(- \frac{\sigma_2^2}{2}\,T + \sigma_2W_t^2\right) \right)\right) - e^{-rT}K,
      \intertext{Les mouvements browniens $(W_t^1)_t$ et $(W_t^2)_t$ sont corrélés de coefficient de correlation $\rho$. On est donc embétés pour calculer cette esperance puisque on ne connait pas la loi jointe de $W_T^1$ et $W_T^2$. On introduit donc un troisième mouvement brownien stantard $(B_t)_t$ définit comme suit $B_t := \frac{W_t^2 - \rho \,W_t^1}{\sqrt{1-\rho^2}}$ indépendant de $W_t^1$. Ainsi,  $W_t^2 = \sqrt{1-\rho^2}B_t + \rho W_t^1$. Maintenant on veut calculer l'esperance d'une fonction de deux mouvement browniens standards \textbf{indépendants}, donc de densité jointe égale au produit des densités.}
      \intertext{}
P^1 &=  \E \left(\max\left(s_0^1 \exp\left(- \frac{\sigma_1^2}{2}\,T + \sigma_1W_T^1\right), s_0^2 \exp\left(- \frac{\sigma_2^2}{2}\,T + \sigma_2(\sqrt{1-\rho^2}B_T + \rho W_T^1)\right) \right)\right)\\
      &\qquad- e^{-rT}K,\\
      &=  \E \Bigg(\exp\left(\sigma_2 \rho W_T^1\right)\,\max\Bigg(s_0^1 \exp\left(- \frac{\sigma_1^2}{2}\,T + \left(\sigma_1 - \sigma_2 \rho\right) W_T^1\right), \\
      &\qquad s_0^2 \exp\left(- \frac{\sigma_2^2}{2}\,T + \sigma_2\sqrt{1-\rho^2}B_T \right) \Bigg)\Bigg) - e^{-rT}K,
        \intertext{Comme $W_t^1$ et $B_t$ sont tout deux des mouvements browniens standards, on peut les écrire comme $W_t^1 \overset{\text{Loi}}{=} \sqrt{t}X$ et $B_t \overset{\text{Loi}}{=} \sqrt{t}Y$ où $X$ et $Y$ sont iid de loi $\N(0,1)$. }
  P^1 &=  \E \Bigg(\exp\left(\sigma_2 \rho \sqrt{T} X\right)\,\max\Bigg(s_0^1 \exp\left(- \frac{\sigma_1^2}{2}\,T + \left(\sigma_1 - \sigma_2 \rho\right) \sqrt{T} X\right), \\
      &\qquad s_0^2 \exp\left(- \frac{\sigma_2^2}{2}\,T + \sigma_2\sqrt{1-\rho^2}\sqrt{T} Y \right) \Bigg)\Bigg) - e^{-rT}K,\\
      &= \int_\R\int_\R \exp\left(\sigma_2 \rho \sqrt{T} x\right)\,\max\Bigg(s_0^1 \exp\left(- \frac{\sigma_1^2}{2}\,T + \left(\sigma_1 - \sigma_2 \rho\right) \sqrt{T} x\right),\\
      &\qquad s_0^2 \exp\left(- \frac{\sigma_2^2}{2}\,T + \sigma_2\sqrt{1-\rho^2}\sqrt{T} y \right) \Bigg) p_X(x) \cdot p_Y(y) \d y \,\d x - e^{-rT}K,\\
  \shortintertext{où $p_X$ et $p_Y$ sont les densités (respectives) des v.a. $X$ et $Y$.}
      &= \int_\R\int_{-\infty}^{A} s^1_0 \exp\left(- \frac{\sigma_1^2}{2}\,T + \sigma_1 \sqrt{T} x\right) p_X(x) \cdot p_Y(y) \d y\, \d x\\
      &\qquad + \int_\R \int_{A}^{+\infty}s^2_0 \exp\left(- \frac{\sigma_2^2}{2}T + \sigma_2\sqrt{1-\rho^2}\sqrt{T} y + \sigma_2 \rho \sqrt{T} x\right) p_X(x) \cdot p_Y(y)\d y\, \d x \\
      &\qquad - e^{-rT}K,
        \shortintertext{où $A = \dfrac{(\sigma_1^2 - \sigma_2^2) + 2\sigma_1}{\rho\sqrt{1-\rho^2}\sqrt{T}} - \sigma_2x$.}        
\end{align*}


\end{document}